\section*{Chapitre \ref{ch:m2:european-and-japanese-government-bonds} --- Obligations d'État européennes et japonaises}

\begin{solution}{exo:m2:european-and-japanese-government-bonds:1}
Espagne $3.633 - 3.18 = 45$, Italie $3.986 - 3.18 = 81$, France $4.00 - 3.18 = 82$
points de base.
\end{solution}

\begin{solution}{exo:m2:european-and-japanese-government-bonds:2}
Le spread s'élargit de $30 - (-10) = 40$ points de base~: 30 viennent de l'Italie, 10
de la baisse de l'Allemagne.
\end{solution}

\begin{solution}{exo:m2:european-and-japanese-government-bonds:3}
\omterm{def:m2:european-and-japanese-government-bonds:ldi}{Gilts} 1\,000, repo 500, coussin 500~: effet de levier 2. Après une baisse de 20~\%~: \omterm{def:m2:european-and-japanese-government-bonds:ldi}{gilts} 800,
repo 500, coussin 300 (60~\% de la valeur initiale), effet de levier $800/300 = 2.67$. L'effet
de levier augmente à mesure que le coussin se réduit~: la perte suivante pèse davantage.
\end{solution}

\begin{solution}{exo:m2:european-and-japanese-government-bonds:4}
Premier ordre~: $0.5/21.8 = 229$ points de base. Deuxième ordre~: la perte $21.8\,x -
294.6\,x^2$ est maximale en $x = 21.8/589 = 3.7\%$ et y vaut $21.8^2/(2
\times 589) = 40.5\%$, en deçà du coussin de 50~\%~: la formule indique que le coussin
n'est jamais épuisé, ce qui est faux. La réévaluation complète donne 351 points de base. Le
développement quadratique sous-estime les pertes loin du point de départ parce que la véritable
courbe de prix y est plus plate que la parabole~: la convexité elle-même diminue à mesure que
les rendements montent.
\end{solution}

\begin{solution}{exo:m2:european-and-japanese-government-bonds:5}
Le problème tenait à un dysfonctionnement sur la partie longue de la courbe, à des ventes forcées des fonds LDI,
et non à l'orientation de la politique monétaire~; la Banque n'a acheté que ce que les fonds
vendaient, soit les \omterm{def:m2:european-and-japanese-government-bonds:ldi}{gilts} conventionnels et indexés longs, pendant le temps
nécessaire aux fonds pour lever des capitaux, puis elle a revendu ces \omterm{def:m2:european-and-japanese-government-bonds:ldi}{gilts}. L'\omterm{def:m2:central-banks-and-the-short-rate:qe}{assouplissement quantitatif}
achète sur toutes les maturités, pendant des années, pour faire baisser les rendements, alors que la
Banque relevait ses taux au même moment et se préparait à réduire son bilan, et une opération
ponctuelle et ciblée a tenu les deux à l'écart.
\end{solution}

\begin{solution}{exo:m2:european-and-japanese-government-bonds:6}
Sur le marché, la banque centrale achète tout ce qui est offert au taux cible. Si le
marché anticipe une hausse ou un abandon de ce taux, détenir des obligations achetées à ce taux
signifie une perte quand le rendement bondit, si bien que les détenteurs vendent d'abord à la
banque centrale~; ses achats et son bilan croissent le plus vite précisément quand la
politique est la moins crédible, et c'est elle qui subit la perte quand le taux cible est abandonné.
\end{solution}

\begin{solution}{exo:m2:european-and-japanese-government-bonds:7}
Après $+100$~: \omterm{def:m2:european-and-japanese-government-bonds:ldi}{gilts} valant 808.4 millions, ventes $2 \times 500 - 808.4 =
\text{GBP}~191.6$ millions. Après $+160$~: \omterm{def:m2:european-and-japanese-government-bonds:ldi}{gilts} 715.6 millions, ventes GBP~284.4
millions. Pour résister à 250 points de base sans vendre, le coussin doit couvrir
la perte à 250~: 39.9~\% de la position (le coussin du fonds, 50~\%, suffit).
\end{solution}

\begin{solution}{exo:m2:european-and-japanese-government-bonds:8}
La hausse des rendements longs réduit la valeur actuelle des engagements d'un régime à prestations définies
plus que celle de ses actions et de ses obligations courtes~: l'illustration propre à la Banque
note que le régime pourrait être globalement mieux loti. La perte
s'est concentrée dans le fonds LDI à effet de levier, dont le coussin s'est érodé, et la
crise relevait de la liquidité~: le collatéral devait être déposé en quelques heures, alors que les
autres actifs des régimes ne pouvaient pas être convertis en liquidités aussi vite.
\end{solution}

\begin{solution}{pb:m2:european-and-japanese-government-bonds:1}
\textbf{1.} Prix plein coupon de 63.41 pour 100 (le \omterm{def:m2:european-and-japanese-government-bonds:ldi}{gilt} représente GBP~1.577 milliard de valeur nominale)~;
repo GBP~500 millions~; coussin GBP~500 millions.
\textbf{2.} \omterm{def:m2:government-bonds:duration}{Duration modifiée} de 21.84~; le DV01 du fonds vaut $10^9 \times 21.84 \times
10^{-4} = \text{GBP}~2.18$ millions par point de base.
\textbf{3.} Effet de levier 2~; 0.44~\% du capital par point de base.
\textbf{4.} 1~\% du coussin vaut GBP~5 millions~: soit environ 2.3 points de base.
\textbf{5.} Pour couvrir davantage ses engagements longs que ne le feraient ses seules
obligations, tout en conservant ses actifs de croissance~: l'effet de levier permet d'acquérir de la duration sans vendre
d'actions.
\textbf{6.} \omterm{def:m2:european-and-japanese-government-bonds:ldi}{Gilts} GBP~808.4 millions, coussin de 308.4 millions (62~\% de la
valeur initiale), effet de levier 2.62.
\textbf{7.} \omterm{def:m2:european-and-japanese-government-bonds:ldi}{Gilts} GBP~715.6 millions, coussin de 215.6 millions (43~\%), effet de levier 3.32.
\textbf{8.} GBP~191.6 millions après 100~; 284.4 millions après 160.
\textbf{9.} Au moins GBP~50 milliards, face à environ 12 milliards de volume quotidien sur
ces maturités, soit plusieurs jours de volume total du marché, écoulés en quelques heures~; le
prix chute jusqu'à ce que d'autres acheteurs apparaissent, ce qui fait monter les rendements et appelle
davantage de ventes.
\textbf{10.} Les fonds LDI détenaient aussi des swaps de taux et d'inflation dans lesquels
ils recevaient le taux fixe~; la hausse des taux a fait perdre de la valeur à ces swaps, et la
marge de variation devait être versée en numéraire chaque jour.
\textbf{11.} 351 points de base.
\textbf{12.} Il vend les \omterm{def:m2:european-and-japanese-government-bonds:ldi}{gilts} qu'il détient en collatéral pour récupérer son prêt,
ce qui accroît l'offre qui a causé la perte.
\textbf{13.} 39.9~\% de la position, GBP~399 millions.
\textbf{14.} Une hausse de 100 points de base coûte 19.2~\% de la position~: un coussin de
environ un cinquième, moins de la moitié de ce qu'exige 250.
\textbf{15.} Le même coussin protège un fonds à courte duration contre des mouvements de rendement bien plus importants
que ne le fait un fonds à longue duration~: la résistance s'exprime en un mouvement de rendement, soit le coussin divisé par
la duration, et non en pourcentage.
\textbf{16.} Parce que la partie longue de la courbe avait cessé de fonctionner~: sans acheteur
la spirale aurait contraint les banques à vendre leur collatéral et perturbé
les marchés de financement. L'opération a été dimensionnée et calibrée dans le temps sur le dysfonctionnement, et non
aux perspectives d'inflation.
\textbf{17.} Gagner du temps pour que les fonds lèvent des capitaux et réduisent leur effet de levier,
sans faire de ces achats un outil de politique monétaire en période de resserrement.
\textbf{18.} Ses engagements ont baissé avec les rendements longs, davantage que ses
actifs hors LDI~: sa situation de financement était, au contraire, meilleure~; ce qui lui manquait était
des liquidités, et rapidement.
\textbf{19.} Résultat nommé~: \emph{rendement d'épuisement} du fonds correspond à une hausse de 351
points de base~; après le mouvement de septembre de 160, son coussin représente 43~\% de la
valeur initiale, et rétablir l'effet de levier exigerait de vendre GBP~284 millions de
\omterm{def:m2:european-and-japanese-government-bonds:ldi}{gilts} dans un marché en baisse.
\textbf{20.} Un effet de levier qu'il fallait compléter en numéraire plus vite que le numéraire
ne pouvait être mobilisé, sur un marché incapable d'absorber les ventes.
\end{solution}

\begin{solution}{iq:m2:european-and-japanese-government-bonds:1}
L'Italie emprunte dans une monnaie qu'elle n'émet pas, si bien qu'elle peut faire défaut en euros~;
l'Allemagne est la référence du crédit le plus sûr et du marché le plus liquide. Le
spread rémunère le risque de crédit italien, l'écart de liquidité et
le \omterm{def:m2:european-and-japanese-government-bonds:spread}{risque de relibellé}, soit la probabilité d'un remboursement dans une monnaie nouvelle et plus faible~; dans une
crise, s'y ajoute la fuite vers le \omterm{def:m2:european-and-japanese-government-bonds:bund}{Bund}, qui fait baisser l'indice de référence.
\iqlookfor{crédit sans planche à billets, liquidité, relibellé.}
\end{solution}

\begin{solution}{iq:m2:european-and-japanese-government-bonds:2}
Après l'annonce budgétaire du gouvernement du 23 septembre, les rendements des \omterm{def:m2:european-and-japanese-government-bonds:ldi}{gilts} longs
ont grimpé plus vite que jamais depuis 2000, de 160 points de base en quelques jours.
Les stratégies LDI des fonds de pension, à effet de levier par le repo et les swaps, ont fait
face à des appels de collatéral qu'ils pouvaient honorer au plus vite en vendant des \omterm{def:m2:european-and-japanese-government-bonds:ldi}{gilts}, ce qui a poussé les rendements
à la hausse~: une spirale. Le 28 septembre, la Banque d'Angleterre a acheté temporairement des \omterm{def:m2:european-and-japanese-government-bonds:ldi}{gilts} longs,
pour 19.3~milliards de livres d'ici au 14 octobre, puis les régulateurs ont exigé que
les fonds LDI puissent résister à un mouvement de 250 points de base.
\iqlookfor{le déclencheur, le mécanisme de levier et le filet de sécurité ciblé.}
\end{solution}

\begin{solution}{iq:m2:european-and-japanese-government-bonds:3}
La banque centrale vise un rendement long et achète tout ce qui est offert à ce niveau~; elle
renonce ainsi à maîtriser son bilan. Sortir du dispositif expose à un risque de saut des
rendements, à mesure que l'acheteur de dernier ressort se retire, à des pertes pour les détenteurs
(et pour la banque centrale sur ses propres avoirs), et à des effets de contagion vers d'autres marchés
dont les investisseurs s'étaient tournés vers l'étranger en quête de rendement. Les sorties se font donc
par étapes, en élargissant la bande avant de mettre fin à la cible.
\iqlookfor{contrôle du prix contre contrôle des quantités et problème de crédibilité.}
\end{solution}

\begin{solution}{iq:m2:european-and-japanese-government-bonds:4}
Décomposer la variation~: variation du rendement de l'émetteur moins variation de celui de l'indice de référence. En crise, l'indice de référence
baisse souvent tandis que l'émetteur monte. Une couverture qui vend à découvert l'indice de référence
contre une position longue sur l'émetteur perd lorsque l'indice de référence remonte~; un
trader positionné sur le spread doit dimensionner les deux jambes par DV01 et savoir laquelle
des deux se déplace.
\iqlookfor{la décomposition et sa conséquence pour la couverture.}
\end{solution}

\begin{solution}{iq:m2:european-and-japanese-government-bonds:5}
OMT (2012)~: achats sans limite ex ante, conditionnés à un programme d'ajustement ou de précaution
avec le Mécanisme européen de stabilité. TPI (2022)~: achats face à des dynamiques de marché
injustifiées et désordonnées qui menacent la transmission de la politique monétaire, conditionnés à
quatre critères de solidité budgétaire et macroéconomique, sans programme. Le premier est un prêteur
en dernier ressort assorti de conditions~; le second est un outil de transmission
soumis à des critères d'éligibilité.
\iqlookfor{la conditionnalité de chacun.}
\end{solution}

\begin{solution}{iq:m2:european-and-japanese-government-bonds:6}
Choisir le choc à partir de l'historique, à la vitesse et sur l'horizon pertinents (de combien les rendements
peuvent-ils varier pendant le temps nécessaire à lever du capital), à partir des cas extrêmes de marchés
comparables, et à partir du stress inversé~: le mouvement qui épuise le capital.
Le test doit capturer la réévaluation complète plutôt que la seule duration, la liquidité (combien peut être vendu,
avec quel impact, dans ce délai), les appels de marge et de collatéral
sur les dérivés, la réaction de la contrepartie (hausse des décotes), et la
rétroaction lorsque de nombreux fonds agissent de la même manière.
\iqlookfor{stress inversé, liquidité et rétroaction, et pas seulement un chiffre plus grand.}
\end{solution}
